\section{Kiểm thử}\label{ch:kiemthu}

\subsection{Chiến lược kiểm thử}
Nhóm kiểm thử theo nhiều mức, mỗi mức nhắm vào một loại rủi ro.

\begin{longtable}{@{}L{2.7cm}L{3.4cm}L{4.7cm}L{3cm}@{}}
\caption{Các mức kiểm thử}\label{tab:mucktest}\\
\toprule
\textbf{Mức} & \textbf{Công cụ, môi trường} & \textbf{Nội dung} & \textbf{Rủi ro nhắm tới} \\ \midrule \endfirsthead
\toprule \textbf{Mức} & \textbf{Công cụ, môi trường} & \textbf{Nội dung} & \textbf{Rủi ro nhắm tới} \\ \midrule \endhead
\bottomrule \endfoot
Unit test API & Vitest, không cần database & Hàm thuần: tính giá, máy trạng thái, công thức nguyên liệu, số ly, chữ ký VNPay, cổng mô phỏng, mã quà, health & Sai công thức, sai quy tắc nghiệp vụ \\
Test HTTP API & Vitest, Supertest, Prisma giả & Toàn bộ pipeline NestJS: guard, \code{ValidationPipe}, mã HTTP, định dạng lỗi & Lỗ hổng phân quyền, dữ liệu sai lọt vào \\
Integration test & Vitest, Supertest, PostgreSQL 16 thật & API kết hợp database: transaction, cập nhật có điều kiện, ràng buộc duy nhất, yêu cầu đồng thời & Bán vượt kho, trừ tiền hai lần, âm điểm \\
Unit test giao diện & Vitest, Testing Library, jsdom & Trang và component: trạng thái tải/rỗng/lỗi, form, giỏ hàng, thanh toán, khôi phục khi lỗi & Giao diện hiển thị sai trạng thái \\
E2E & Playwright, Chromium, API, web, PostgreSQL thật đã seed & Luồng của người dùng trên trình duyệt thật & Các phần riêng lẻ đúng nhưng ghép lại sai \\
Phi chức năng & Script Node.js trong CI & Độ trễ \code{GET /products} (NFR-01), quét secret trong repository (NFR-03) & API chậm dần, lộ secret khi commit \\
Kiểm thử thủ công trên production & Trình duyệt máy tính và điện thoại & Toàn bộ luồng khách và nhân viên, 4 phương thức thanh toán, VNPay thật (sandbox) & Lỗi chỉ xuất hiện trên môi trường thật \\
\end{longtable}

\textbf{Nguyên tắc: các tình huống đồng thời phải test trên PostgreSQL thật.} Cập nhật có điều kiện và ràng buộc duy nhất chỉ có ý nghĩa khi có transaction của database. Mock Prisma sẽ không phát hiện được lỗi bán vượt kho. Vì vậy integration test chạy trong CI với PostgreSQL trong service container, không bỏ qua khi thiếu database.

\subsection{Số lượng và kết quả}
Bảng \ref{tab:sotest} là số test tại \code{v1.1.1}. Nhóm chạy lại toàn bộ test (trừ E2E) trên mã của \code{main} lúc 04:24 ngày 17/11/2026: unit test API 62/62, test HTTP API 16/16, integration test (6 tệp, PostgreSQL 16) 53/53, unit test web 72/72.

E2E: 8/8 test đạt khi chạy trên máy với mã của PR \#55. CI của PR \#55 và các PR sau đều xanh ở cả hai job.

\begin{table}[H]
\centering
\caption{Số test theo phiên bản}\label{tab:sotest}
\begin{tabular}{@{}L{5cm}rrL{5cm}@{}}
\toprule
\textbf{Loại} & \textbf{\code{v1.0.0}} & \textbf{\code{v1.1.1}} & \textbf{Lệnh} \\ \midrule
Unit API & 55 & 62 & \code{npm test -w apps/\allowbreak{}api} \\
Test HTTP API (Prisma giả) & 16 & 16 & \code{npm run test:e2e -w apps/\allowbreak{}api} \\
Integration API (PostgreSQL thật) & 44 & 53 & \code{npm run test:int -w apps/\allowbreak{}api} \\
Unit web & 55 & 72 & \code{npm test -w apps/\allowbreak{}web} \\
E2E trình duyệt (Playwright) & 0 & 8 & \code{cd tests/\allowbreak{}e2e \&\& npx playwright test} \\
\midrule
\textbf{Tổng} & \textbf{170} & \textbf{211} & \\
\midrule
Phi chức năng (sau \code{v1.1.1}) & -- & 6 & mục \ref{sec:nfrtest} \\
\bottomrule
\end{tabular}
\end{table}

\subsection{Danh sách test case}
Mã test (bảng \ref{tab:tvkyhieu}) đánh theo thứ tự \code{it(...)} trong tệp và được dùng trong ma trận truy vết. Tên test case lấy nguyên từ mã nguồn. Nhiều test case ghi luôn mã FR mà nó kiểm chứng.

{\small
\begin{longtable}{@{}L{2.6cm}L{5.6cm}L{5.7cm}@{}}
\caption{Ký hiệu mã test dùng trong ma trận truy vết}\label{tab:tvkyhieu}\\
\toprule
\textbf{Mã} & \textbf{Loại} & \textbf{Tệp} \\ \midrule \endfirsthead
\toprule \textbf{Mã} & \textbf{Loại} & \textbf{Tệp} \\ \midrule \endhead
\bottomrule \endfoot
UT-* & Unit test API & \code{apps/\allowbreak{}api/\allowbreak{}src/\allowbreak{}**/\allowbreak{}*.\allowbreak{}spec.\allowbreak{}ts} \\
WEB-* & Unit/component test giao diện & \code{apps/\allowbreak{}web/\allowbreak{}**/\allowbreak{}*.\allowbreak{}test.\allowbreak{}tsx} \\
API-T01--T07 & Test HTTP API (Prisma giả) & \code{apps/\allowbreak{}api/\allowbreak{}test/\allowbreak{}app.\allowbreak{}e2e-spec.\allowbreak{}ts} (theo thứ tự trong tệp) \\
AUTH-T01--T09 & Test HTTP xác thực, phân quyền, validate & \code{apps/\allowbreak{}api/\allowbreak{}test/\allowbreak{}auth-guard.\allowbreak{}e2e-spec.\allowbreak{}ts} \\
INT-A01--A11 & Integration: xác thực, đơn hàng & \code{apps/\allowbreak{}api/\allowbreak{}test/\allowbreak{}integration/\allowbreak{}auth-orders.\allowbreak{}int-spec.\allowbreak{}ts} \\
INT-P01--P10 & Integration: thanh toán Ví/Thẻ, idempotency, tồn kho đồng thời & \code{apps/\allowbreak{}api/\allowbreak{}test/\allowbreak{}integration/\allowbreak{}payments.\allowbreak{}int-spec.\allowbreak{}ts} \\
INT-Q01--Q09 & Integration: thanh toán QR & \code{apps/\allowbreak{}api/\allowbreak{}test/\allowbreak{}integration/\allowbreak{}qr-payments.\allowbreak{}int-spec.\allowbreak{}ts} \\
INT-V01--V06 & Integration: VNPay sandbox & \code{apps/\allowbreak{}api/\allowbreak{}test/\allowbreak{}integration/\allowbreak{}vnpay.\allowbreak{}int-spec.\allowbreak{}ts} \\
INT-R01--R09 & Integration: số ly làm được, công thức & \code{apps/\allowbreak{}api/\allowbreak{}test/\allowbreak{}integration/\allowbreak{}recipes-availability.\allowbreak{}int-spec.\allowbreak{}ts} \\
INT-S01--S08 & Integration: kho, đổi quà, tự hủy đơn & \code{apps/\allowbreak{}api/\allowbreak{}test/\allowbreak{}integration/\allowbreak{}staff-rewards.\allowbreak{}int-spec.\allowbreak{}ts} \\
E2E-T01--T07 & E2E trình duyệt (Playwright) trên API, web, PostgreSQL thật & \code{tests/\allowbreak{}e2e/\allowbreak{}brewlite-flow.\allowbreak{}spec.\allowbreak{}ts} \\
\end{longtable}
}

\subsubsection{Test HTTP API}
{\small
\begin{longtable}{@{}L{1.9cm}L{10.6cm}L{1.3cm}@{}}
\caption{Test HTTP API với Prisma giả (\code{app.\allowbreak{}e2e-spec.\allowbreak{}ts})}\label{tab:apit}\\
\toprule
\textbf{Mã} & \textbf{Test case (tên trong mã nguồn)} & \textbf{Kết quả} \\ \midrule \endfirsthead
\toprule \textbf{Mã} & \textbf{Test case} & \textbf{Kết quả} \\ \midrule \endhead
\bottomrule \endfoot
API-T01 & GET /health trả 200, status ok, db up & Đạt \\
API-T02 & GET /health trả 503 khi database lỗi & Đạt \\
API-T03 & GET /products trả danh sách menu & Đạt \\
API-T04 & GET /products/:id trả chi tiết kèm size và mức đường, đá & Đạt \\
API-T05 & GET /products/:id trả 404 khi không có sản phẩm & Đạt \\
API-T06 & GET /products/:id trả 400 khi id không phải số & Đạt \\
API-T07 & trả 404 cho đường dẫn không tồn tại & Đạt \\
\end{longtable}
}
{\small
\begin{longtable}{@{}L{1.9cm}L{10.6cm}L{1.3cm}@{}}
\caption{Test HTTP xác thực, phân quyền, validate (\code{auth-guard.\allowbreak{}e2e-spec.\allowbreak{}ts})}\label{tab:autht}\\
\toprule
\textbf{Mã} & \textbf{Test case (tên trong mã nguồn)} & \textbf{Kết quả} \\ \midrule \endfirsthead
\toprule \textbf{Mã} & \textbf{Test case} & \textbf{Kết quả} \\ \midrule \endhead
\bottomrule \endfoot
AUTH-T01 & POST /orders không có JWT trả 401 & Đạt \\
AUTH-T02 & JWT ký bằng secret khác trả 401 & Đạt \\
AUTH-T03 & CUSTOMER đổi trạng thái đơn nhận 403 & Đạt \\
AUTH-T04 & CUSTOMER xem danh sách đơn của nhân viên nhận 403 & Đạt \\
AUTH-T05 & STAFF đặt trạng thái không cho phép (PAID) nhận 400 & Đạt \\
AUTH-T06 & đặt đơn với dữ liệu sai hoặc trường lạ trả 400 & Đạt \\
AUTH-T07 & POST /payments thiếu Idempotency-Key trả 400 & Đạt \\
AUTH-T08 & POST /payments bằng thẻ mà thiếu số thẻ trả 400 & Đạt \\
AUTH-T09 & đăng ký với email sai hoặc mật khẩu ngắn trả 400 & Đạt \\
\end{longtable}
}

\subsubsection{Integration test}
{\small
\begin{longtable}{@{}L{1.9cm}L{10.6cm}L{1.3cm}@{}}
\caption{Integration test xác thực và đơn hàng (\code{auth-orders.\allowbreak{}int-spec.\allowbreak{}ts})}\label{tab:inta}\\
\toprule
\textbf{Mã} & \textbf{Test case (tên trong mã nguồn)} & \textbf{Kết quả} \\ \midrule \endfirsthead
\toprule \textbf{Mã} & \textbf{Test case} & \textbf{Kết quả} \\ \midrule \endhead
\bottomrule \endfoot
INT-A01 & đăng ký, mật khẩu được băm, đăng nhập nhận JWT, /auth/me trả đúng người & Đạt \\
INT-A02 & email trùng trả 409, sai mật khẩu trả 401 & Đạt \\
INT-A03 & tính giá ở máy chủ, áp mã giảm giá, chưa trừ kho & Đạt \\
INT-A04 & mã hết hạn, chưa đủ đơn tối thiểu, hết lượt bị từ chối kèm lý do & Đạt \\
INT-A05 & thiếu nguyên liệu trả 400, báo món thiếu và không tạo đơn & Đạt \\
INT-A06 & sản phẩm hoặc topping không tồn tại trả 400 & Đạt \\
INT-A07 & GET /orders/me chỉ trả đơn của người đăng nhập, mới nhất trước & Đạt \\
INT-A08 & người khác xem đơn nhận 404, STAFF xem được & Đạt \\
INT-A09 & STAFF lọc đơn theo trạng thái; CUSTOMER bị 403 & Đạt \\
INT-A10 & chuyển sai trạng thái trả 400 và đơn giữ nguyên & Đạt \\
INT-A11 & hủy đơn chưa thanh toán trả lại lượt dùng mã khuyến mãi, không đụng kho & Đạt \\
\end{longtable}
}
{\small
\begin{longtable}{@{}L{1.9cm}L{10.6cm}L{1.3cm}@{}}
\caption{Integration test thanh toán Ví/Thẻ, idempotency, tồn kho đồng thời (\code{payments.\allowbreak{}int-spec.\allowbreak{}ts})}\label{tab:intp}\\
\toprule
\textbf{Mã} & \textbf{Test case (tên trong mã nguồn)} & \textbf{Kết quả} \\ \midrule \endfirsthead
\toprule \textbf{Mã} & \textbf{Test case} & \textbf{Kết quả} \\ \midrule \endhead
\bottomrule \endfoot
INT-P01 & FR-08, FR-12, FR-14: thanh toán Ví thành công, trừ đúng kho, cộng điểm & Đạt \\
INT-P02 & FR-09: gửi lại cùng Idempotency-Key trả cùng kết quả, chỉ một bản ghi Payment & Đạt \\
INT-P03 & FR-09: cùng khóa nhưng nội dung khác trả 409; đơn đã PAID không thanh toán lại & Đạt \\
INT-P04 & FR-08: thẻ bị từ chối $\rightarrow$ PAYMENT\_FAILED, hoàn kho; thử lại bằng khóa mới thì PAID & Đạt \\
INT-P05 & không thanh toán được đơn của người khác (404) hoặc thiếu khóa (400) & Đạt \\
INT-P06 & FR-12: kho vừa hết lúc thanh toán $\rightarrow$ không trừ tiền, không trừ kho, PAYMENT\_FAILED kèm lý do & Đạt \\
INT-P07 & FR-12: N khách thanh toán đồng thời không bán vượt kho, tồn kho không âm & Đạt \\
INT-P08 & FR-09: hai yêu cầu đồng thời cho cùng một đơn chỉ thanh toán một lần & Đạt \\
INT-P09 & FR-09: bấm thanh toán hai lần liên tiếp cùng khóa chỉ tạo một Payment & Đạt \\
INT-P10 & FR-10, FR-12, FR-14: nhân viên xử lý đơn; hủy đơn đã PAID hoàn kho và trừ lại điểm & Đạt \\
\end{longtable}
}

Hai test quan trọng nhất của bảng \ref{tab:intp}:
\begin{itemize}
\item \textbf{INT-P07:} nhiều khách thanh toán đồng thời khi kho chỉ đủ cho một phần. Số đơn PAID đúng bằng số ly kho cho phép, tồn kho không âm, các đơn còn lại PAYMENT\_FAILED và không bị trừ tiền.
\item \textbf{INT-P08:} hai yêu cầu đồng thời cho cùng một đơn. Chỉ một lần thanh toán được ghi nhận.
\end{itemize}

{\small
\begin{longtable}{@{}L{1.9cm}L{10.6cm}L{1.3cm}@{}}
\caption{Integration test thanh toán QR (\code{qr-payments.\allowbreak{}int-spec.\allowbreak{}ts})}\label{tab:intq}\\
\toprule
\textbf{Mã} & \textbf{Test case (tên trong mã nguồn)} & \textbf{Kết quả} \\ \midrule \endfirsthead
\toprule \textbf{Mã} & \textbf{Test case} & \textbf{Kết quả} \\ \midrule \endhead
\bottomrule \endfoot
INT-Q01 & tạo mã QR: token ngẫu nhiên, link tới trang xác nhận, ảnh QR; tạo lại khi còn hạn thì trả cùng mã & Đạt \\
INT-Q02 & trang điện thoại xem được số tiền và tài khoản thử, không lộ số dư; token sai trả 404 & Đạt \\
INT-Q03 & sai PIN: 400, không trừ tiền, không đụng kho, cho nhập lại & Đạt \\
INT-Q04 & đúng PIN, đủ số dư: đơn PAID, trừ đúng số tiền, trừ kho, cộng điểm; khách thấy trạng thái PAID & Đạt \\
INT-Q05 & xác nhận lặp lại chỉ thành công một lần (tuần tự và đồng thời) & Đạt \\
INT-Q06 & lần xác nhận đến sau trong lúc lần trước vừa ghi xong (kết quả được trả lại) vẫn nhận 409, không trừ tiền lần hai & Đạt \\
INT-Q07 & thiếu số dư: không trừ tiền, hoàn kho, đơn PAYMENT\_FAILED; tạo mã mới và trả bằng tài khoản khác thì PAID & Đạt \\
INT-Q08 & QR hết hạn: không trừ tiền, trạng thái EXPIRED & Đạt \\
INT-Q09 & chỉ chủ đơn tạo được QR; chưa đăng nhập trả 401; đơn đã thanh toán trả 409 & Đạt \\
\end{longtable}
}
{\small
\begin{longtable}{@{}L{1.9cm}L{10.6cm}L{1.3cm}@{}}
\caption{Integration test VNPay sandbox (\code{vnpay.\allowbreak{}int-spec.\allowbreak{}ts})}\label{tab:intv}\\
\toprule
\textbf{Mã} & \textbf{Test case (tên trong mã nguồn)} & \textbf{Kết quả} \\ \midrule \endfirsthead
\toprule \textbf{Mã} & \textbf{Test case} & \textbf{Kết quả} \\ \midrule \endhead
\bottomrule \endfoot
INT-V01 & tạo URL thanh toán đã ký đúng: số tiền $\times$100, mã website, IP khách, chữ ký kiểm tra lại được & Đạt \\
INT-V02 & Return URL hợp lệ, mã 00: đơn PAID, trừ kho, cộng điểm; gọi lại và IPN không tạo thêm thanh toán & Đạt \\
INT-V03 & sai chữ ký: từ chối, đơn giữ nguyên PENDING, IPN trả 97 & Đạt \\
INT-V04 & chữ ký đúng nhưng số tiền khác đơn: từ chối, đơn giữ nguyên, IPN trả 04 & Đạt \\
INT-V05 & khách hủy ở VNPay (mã 24): PAYMENT\_FAILED kèm lý do, hoàn kho, cho thanh toán lại bằng giao dịch mới & Đạt \\
INT-V06 & giao dịch không tồn tại: Return 404, IPN 01 & Đạt \\
\end{longtable}
}
{\small
\begin{longtable}{@{}L{1.9cm}L{10.6cm}L{1.3cm}@{}}
\caption{Integration test số ly làm được và công thức (\code{recipes-availability.\allowbreak{}int-spec.\allowbreak{}ts})}\label{tab:intr}\\
\toprule
\textbf{Mã} & \textbf{Test case (tên trong mã nguồn)} & \textbf{Kết quả} \\ \midrule \endfirsthead
\toprule \textbf{Mã} & \textbf{Test case} & \textbf{Kết quả} \\ \midrule \endhead
\bottomrule \endfoot
INT-R01 & size lớn dùng nhiều nguyên liệu hơn: cà phê 60g $\rightarrow$ S 3 ly, M 2 ly, L 2 ly; size L không đủ dù size S còn đủ & Đạt \\
INT-R02 & hai dòng cùng món tranh nhau nguyên liệu: kho đủ 5 ly, dòng 1 đặt 3 ly thì dòng 2 còn tối đa 2 ly & Đạt \\
INT-R03 & topping trừ nguyên liệu riêng: trân châu 80g chỉ đủ 2 ly có topping; hết trân châu thì ly có topping báo hết, ly không topping vẫn đặt được & Đạt \\
INT-R04 & giới hạn 50 ly mỗi dòng dù nguyên liệu còn nhiều; không trả số tồn kho & Đạt \\
INT-R05 & dữ liệu sai: giỏ rỗng, số lượng 0, món hoặc topping không tồn tại, quá 5 topping $\rightarrow$ 400 & Đạt \\
INT-R06 & khớp với lúc đặt đơn: availability báo không đủ thì POST /orders cũng từ chối; báo đủ thì đặt được & Đạt \\
INT-R07 & chỉ STAFF xem được: CUSTOMER 403, chưa đăng nhập 401 & Đạt \\
INT-R08 & công thức theo size (làm tròn lên), số ly S/M/L theo tồn kho thật, nguyên liệu giới hạn và số phần topping & Đạt \\
INT-R09 & số ly cập nhật theo kho: thanh toán thì giảm, nhập thêm nguyên liệu thì tăng ngay & Đạt \\
\end{longtable}
}
{\small
\begin{longtable}{@{}L{1.9cm}L{10.6cm}L{1.3cm}@{}}
\caption{Integration test kho, đổi quà, tự hủy đơn (\code{staff-rewards.\allowbreak{}int-spec.\allowbreak{}ts})}\label{tab:ints}\\
\toprule
\textbf{Mã} & \textbf{Test case (tên trong mã nguồn)} & \textbf{Kết quả} \\ \midrule \endfirsthead
\toprule \textbf{Mã} & \textbf{Test case} & \textbf{Kết quả} \\ \midrule \endhead
\bottomrule \endfoot
INT-S01 & STAFF xem tồn kho kèm cờ OUT, LOW, OK; CUSTOMER 403; chưa đăng nhập 401 & Đạt \\
INT-S02 & nhập thêm cộng dồn đúng; số lượng không dương trả 400; nguyên liệu không tồn tại trả 404 & Đạt \\
INT-S03 & sau khi nhập thêm, món hết hàng bán lại được & Đạt \\
INT-S04 & GET /rewards chỉ trả quà đang mở, xếp theo điểm & Đạt \\
INT-S05 & đủ điểm: trừ điểm, trả mã quà, ghi giao dịch REDEEM; không đủ điểm: 400, điểm giữ nguyên & Đạt \\
INT-S06 & quà ngừng hoặc không tồn tại trả 404; chưa đăng nhập trả 401 & Đạt \\
INT-S07 & đổi đồng thời không làm âm điểm: 100 điểm, 5 yêu cầu 50 điểm $\rightarrow$ đúng 2 thành công & Đạt \\
INT-S08 & hủy đơn PENDING quá 15 phút, trả lượt mã khuyến mãi, không đụng đơn mới hoặc đơn đã thanh toán, không đụng kho & Đạt \\
\end{longtable}
}

\subsubsection{Test E2E trên trình duyệt}\label{sec:e2e}
Các kịch bản E2E chạy trên dữ liệu seed và tạo tài khoản mới mỗi lần chạy. Test thay đổi dữ liệu (tạo đơn, nhập kho), nên chỉ chạy trên database riêng, không chạy trên production.

{\small
\begin{longtable}{@{}L{1.9cm}L{10.6cm}L{1.3cm}@{}}
\caption{Test E2E trên trình duyệt (\code{tests/\allowbreak{}e2e/\allowbreak{}brewlite-flow.\allowbreak{}spec.\allowbreak{}ts})}\label{tab:e2e}\\
\toprule
\textbf{Mã} & \textbf{Kịch bản} & \textbf{Kết quả} \\ \midrule \endfirsthead
\toprule \textbf{Mã} & \textbf{Kịch bản} & \textbf{Kết quả} \\ \midrule \endhead
\bottomrule \endfoot
E2E-T01 & menu hiển thị món từ API và mở trang chi tiết & Đạt \\
E2E-T02 & đặt món nhiều topping, thanh toán thất bại rồi thành công, theo dõi đơn và lịch sử & Đạt \\
E2E-T03 & chưa đăng nhập bị chuyển sang /login, giỏ hàng vẫn giữ & Đạt \\
E2E-T04 & thanh toán QR bằng tài khoản thử, trang thanh toán tự chuyển sang thành công & Đạt \\
E2E-T05 & nhân viên xem công thức, nhập kho làm số ly tăng & Đạt \\
E2E-T06a & trang chủ, menu, giỏ hàng không cuộn ngang trên màn hình 375px & Đạt \\
E2E-T06b & manifest hợp lệ và biểu tượng tải được & Đạt \\
E2E-T07 & availability giới hạn 50 ly mỗi dòng và không trả tồn kho & Đạt \\
\end{longtable}
}

\subsection{Kiểm thử phi chức năng}\label{sec:nfrtest}
Sau \code{v1.1.1}, nhóm bổ sung test tự động cho hai NFR trước đó chỉ kiểm chứng thủ công. Đây là script Node.js không cần thư viện ngoài, chạy trong CI.

{\small
\begin{longtable}{@{}L{1.9cm}L{1.3cm}L{7.6cm}L{3cm}@{}}
\caption{Test phi chức năng}\label{tab:nfrtest}\\
\toprule
\textbf{Mã} & \textbf{Yêu cầu} & \textbf{Cách kiểm tra} & \textbf{Vị trí trong CI} \\ \midrule \endfirsthead
\toprule \textbf{Mã} & \textbf{Yêu cầu} & \textbf{Cách kiểm tra} & \textbf{Vị trí trong CI} \\ \midrule \endhead
\bottomrule \endfoot
PERF-T01 & NFR-01 & 1 yêu cầu khởi động, rồi 20 yêu cầu \code{GET /products} liên tiếp trên dữ liệu seed; tính min, trung bình, p95 (nearest-rank), max; trượt khi p95 $\ge$ 500 ms, API lỗi hoặc danh sách rỗng & Job E2E, sau khi API và web đã khởi động \\
PERF-U01--U04 & NFR-01 & Unit test hàm phân vị: p95 của 20 mẫu là mẫu thứ 19; 1 yêu cầu chậm trong 20 vẫn đạt, 2 yêu cầu chậm thì trượt & Job chính, sau unit test \\
SEC-T01 & NFR-03 & Quét các tệp Git theo dõi: tệp \code{.env}, khóa riêng, token GitHub, khóa AWS, URL deploy hook, chuỗi kết nối database có mật khẩu, \code{JWT\_SECRET} hoặc \code{VNPAY\_\allowbreak{}HASH\_\allowbreak{}SECRET} có giá trị thật; che giá trị khi in & Job chính, ngay sau \code{npm ci} \\
\end{longtable}}

\textbf{Kết quả} chạy trên máy ngày 17/11/2026 (API build từ nhánh \code{test/\allowbreak{}nfr-perf}, PostgreSQL 16, dữ liệu seed): PERF-T01 đạt, 20/20 yêu cầu dưới 500 ms (min 2,9; trung bình 5,2; p95 9,5; max 10,7 ms); đặt ngưỡng 1 ms thì script báo ``KHÔNG ĐẠT'' và trả mã thoát 1. PERF-U01--U04 đạt 4/4. SEC-T01 quét 339 tệp, không phát hiện secret; trên repository thử có chèn \code{.env}, chuỗi kết nối database ngoài, \code{JWT\_SECRET} ngẫu nhiên và URL deploy hook, script bắt đủ 4 vị trí và che giá trị.

\textbf{Giới hạn:} PERF-T01 đo độ trễ tuần tự từ một máy, chưa phải load test đồng thời; trong CI, API và database chạy cùng runner nên số đo thấp hơn production. SEC-T01 quét theo mẫu, không quét lịch sử commit.

\subsection{Kiểm thử thủ công trên production}
Ngoài test tự động, nhóm kiểm thử thủ công trên production ở hai đợt:
\begin{itemize}
\item \textbf{16/11/2026:} VNPay sandbox, đối soát 2 giao dịch với Merchant Admin.
\item \textbf{17/11/2026:} toàn bộ luồng trên \code{v1.1.0}, sau đó kiểm tra lại sau khi phát hành \code{v1.1.1}.
\end{itemize}

Các trường hợp được thử, kèm ảnh bằng chứng trong Phụ lục \ref{app:bangchung}:

\begin{longtable}{@{}L{4.6cm}L{6.2cm}L{1.6cm}L{1.4cm}@{}}
\caption{Kiểm thử thủ công trên production}\label{tab:thucong}\\
\toprule
\textbf{Trường hợp} & \textbf{Kết quả quan sát} & \textbf{Ảnh} & \textbf{Kết quả} \\ \midrule \endfirsthead
\toprule \textbf{Trường hợp} & \textbf{Kết quả quan sát} & \textbf{Ảnh} & \textbf{Kết quả} \\ \midrule \endhead
\bottomrule \endfoot
Chưa đăng nhập bấm thanh toán & Chuyển sang \code{/\allowbreak{}login?next=/\allowbreak{}checkout}, giỏ hàng giữ nguyên & 34 & Đạt \\
Thanh toán Ví với mã \code{GIAM10} & Đơn \#12 PAID; $1.426.000 - 142.600 = 1.283.400$ đồng; cộng 128 điểm & 35--37 & Đạt \\
Thẻ bị từ chối rồi thử lại & Báo ``Thẻ bị ngân hàng từ chối'', có nút tạo đơn mới; thẻ hợp lệ thì đơn \#13 PAID, $186.000 \times 0{,}9 = 167.400$ đồng & 38--39 & Đạt \\
QR thành công & Điện thoại xác nhận đơn \#19; trang thanh toán trên máy tính tự chuyển sang thành công & 40--42 & Đạt \\
QR thiếu số dư & Đơn \#14 (390.600 đồng) bằng TEST-0001 sau khi tài khoản đã trả 882.000 đồng cho đơn \#7: báo ``Số dư tài khoản không đủ'', không trừ tiền & 43--46 & Đạt \\
VNPay sandbox & Đơn \#16, 522.000 đồng qua thẻ NCB; trở về web báo thành công, cộng 52 điểm & 47--48 & Đạt \\
Nhân viên chuyển trạng thái & Đơn \#2 qua Chờ pha $\rightarrow$ Đang pha $\rightarrow$ Sẵn sàng $\rightarrow$ Hoàn tất; khách thấy ``Hoàn tất'' & 52--56 & Đạt \\
Đơn quá hạn tự hủy & Đơn \#11 ``Quá hạn thanh toán, đơn tự hủy'' & 57 & Đạt \\
Nhập kho cập nhật số ly & Nhập 200 g đá viên: 116/93/77 $\rightarrow$ 118/94/78 & 59--61 & Đạt \\
Giao diện điện thoại & Header, thực đơn, thẻ món hiển thị đủ, không cuộn ngang & 62 & Đạt \\
Biểu tượng PWA & \code{v1.1.0}: Chrome báo lỗi biểu tượng; \code{v1.1.1}: hiện biểu tượng PNG 192 và 512 & 63, 65 & Lỗi, đã sửa \\
Nhân viên đăng xuất & \code{v1.1.0}: ở lại form đăng nhập nhân viên; \code{v1.1.1}: về trang chủ & 66--67 & Lỗi, đã sửa \\
\end{longtable}
